\scanpage{60}
2°) $\pi^\Gamma=(1)$ (i\unsure{.|:}e there are no other
fixed poin\textsuperscript{t}s in $\widetilde U$ over the $\unsure{=|\overline m}$ $x\in U(K)$,
and finally

3°) there is not at the $\unsure{=|\overline m}$ time a fixed
point $\deleted{\unsure{x|\mathrm{at}}}\ \widetilde x$ at infinity --- i.e there is
no $L_i$ normalized by $\Gamma$, i.e. a splitting
of the first type is not at the $\unsure{=|\overline m}$ time
of the second (a fact which we have
established directly, previously).

\margin{This is a special case of a ``conjectural consequence'' of the conjectures of $\S1$ --- except that one should $\unsure{=|\overline m}$ write $\pi^{\Gamma^\circ}=(1)$}

On the other hand, in the case of fixed points
at $\infty$, uniqueness of the image in $X$
means that one and the $\unsure{=|\overline m}$ effective action cannot
normalize both an $L_i$ and an $L_j$
($i\ne j$) --- a fact also established except
in the case $(g,\nu)=(0,3)$ --- and uniqueness
--- over a fixed $i\in\bar I$ means
that the $L_i$ \inserted{($i$ fixed)} normalized by $\Gamma$ is unique,
which is a weakening of the
relation $L_i=\unsure{\deleted{\unsure{Z}}\mathrm{er}|\deleted{\unsure{Z}}\mathrm{en}|\deleted{\unsure{Z}}_{\mathrm{en}}|\deleted{\unsure{Z}}_{\pi}}\pi^{\Gamma^\circ}$
for these actions, conjectured above.
\margin{59}
